\section{Descripción del dataset}
\label{cap:dataset}

\subsection{Origen}

ListenBrainz publica sus datos de escuchas en varios formatos \citep{listenbrainzdumps}: volcados completos (del orden de cientos de GB), volcados incrementales con las escuchas recibidas desde el incremental anterior y exportaciones orientadas a Spark. Elegimos un incremental porque cumple con holgura el mínimo de 50\,000 registros de la consigna, se descarga y procesa en una laptop y, al ser parte de una serie que se publica periódicamente, encaja con la idea de actualizaciones incrementales. Para que todo el grupo trabaje con la misma versión fijamos un incremental concreto \citep{incremental2746}. La tabla~\ref{tab:dataset} resume sus características.

\begin{table}[H]
\centering
\caption{Volcado utilizado en el proyecto}
\label{tab:dataset}
\begin{tabularx}{\textwidth}{l Y}
\toprule
Propiedad & Valor \\
\midrule
Origen & MetaBrainz, directorio de incrementales de ListenBrainz \\
Corte & 27 de octubre de 2026, volcado número 2746 \\
Tipo & \code{listenbrainz-listens-dump} \\
Licencia & CC0 \\
Archivo comprimido & \code{.tar.zst}, aproximadamente 238~MB \\
Archivo extraído & JSON Lines, aproximadamente 2,6~GB \\
Registros & 4\,855\,809 \\
Ruta en el proyecto & \code{data/processed/listens_full.jsonl} \\
\bottomrule
\end{tabularx}
\end{table}

Un incremental no es todo ListenBrainz: contiene los eventos enviados a la plataforma desde el incremental anterior. Para reconstruir una réplica completa haría falta un volcado completo y aplicar sobre él todos los incrementales, incluidas las eliminaciones. En este proyecto estudiamos expresamente ese lote de ingesta.

La descarga la hacen \code{scripts/download_dataset.sh} (Linux y macOS) y \code{scripts/download_dataset.ps1} (Windows). Ambos bajan primero el archivo \code{.sha256} oficial, descargan el comprimido, comprueban que la huella SHA-256 coincida y recién entonces extraen el archivo de escuchas. La extracción escribe a un archivo temporal \code{.part} y lo renombra al terminar, así nunca queda un archivo incompleto con el nombre final. Si el comprimido ya existe y su huella es correcta, no se vuelve a descargar. MetaBrainz solo mantiene publicados los incrementales de las últimas semanas y el 2746 ya no está en su servidor, así que los scripts aceptan el comprimido y el \code{.sha256} guardados en \code{data/raw/}: si el checksum oficial no se puede descargar, verifican el archivo con la copia local.

\subsection{Estructura de cada registro}

Cada línea del archivo es un objeto JSON que representa una escucha, con la estructura que documenta ListenBrainz para el volcado de escuchas \citep{listenbrainzdumps}. El siguiente ejemplo tiene la estructura real pero datos ficticios, para no publicar datos personales:

\begin{lstlisting}[language=JSON]
{
  "user_id": 12345,
  "user_name": "usuario_ejemplo",
  "timestamp": 1793061000,
  "recording_msid": "identificador-ejemplo",
  "track_metadata": {
    "track_name": "Cancion de ejemplo",
    "artist_name": "Artista de ejemplo",
    "release_name": "Album de ejemplo",
    "additional_info": {}
  }
}
\end{lstlisting}

La tabla~\ref{tab:campos} describe cada campo.

\begin{table}[H]
\centering
\caption{Campos de una escucha}
\label{tab:campos}
\begin{tabularx}{\textwidth}{l l Y}
\toprule
Campo & Significado & Observaciones \\
\midrule
\code{user_id} & Identificador interno del usuario & Entero en todo el volcado \\
\code{user_name} & Nombre de usuario & Dato público; se evita mostrarlo sin necesidad \\
\code{timestamp} & Fecha de la escucha, Unix UTC & En este volcado el campo se llama \code{timestamp}, no \code{listened_at} como en la API \\
\code{recording_msid} & Identificador de la grabación & Presente en todo el volcado \\
\code{track_name} & Nombre de la canción & Texto libre \\
\code{artist_name} & Nombre del artista & Texto libre \\
\code{release_name} & Álbum o lanzamiento & Opcional \\
\code{additional_info} & Metadatos adicionales & Opcional y variable según el cliente \\
\bottomrule
\end{tabularx}
\end{table}

Los cuatro últimos campos están anidados dentro de \code{track_metadata}. Conviene no confundir \code{timestamp} con \code{listened_at}: algunas respuestas de la API de ListenBrainz usan \code{listened_at} \citep{listenbrainzjson}, pero el archivo descargado usa \code{timestamp}, y todos los scripts del repositorio están adaptados a ese formato.

\subsection{Volumen y exploración}

La exploración inicial del archivo dio los valores de la tabla~\ref{tab:exploracion}.

\begin{table}[H]
\centering
\caption{Exploración del incremental 2746}
\label{tab:exploracion}
\begin{tabular}{l r}
\toprule
Medida & Valor \\
\midrule
Escuchas & 4\,855\,809 \\
Usuarios únicos (\code{user_id}) & 26\,394 \\
Nombres de artista distintos & 286\,291 \\
Combinaciones artista y canción & 1\,066\,683 \\
\code{recording_msid} distintos & 1\,420\,842 \\
Primera fecha de escucha (UTC) & 2005-02-13 \\
Última fecha de escucha (UTC) & 2026-10-27 \\
\bottomrule
\end{tabular}
\end{table}

Las cifras describen este incremento, no a la población de ListenBrainz. Los conteos por nombre tampoco equivalen a artistas o canciones únicos: hay homónimos, grafías distintas del mismo artista y registros sin identificador musical. Por eso los 1\,420\,842 \code{recording_msid} distintos no coinciden con las 1\,066\,683 combinaciones de artista y título: son dos criterios de identidad diferentes. El grafo de Neo4j usa el identificador y no fusiona canciones por título.

El rango de fechas muestra otra particularidad: un registro enviado en 2026 puede tener fecha de escucha de 2005, porque los usuarios suben historiales antiguos. La ingesta, en cambio, se concentra en los últimos días de octubre de 2026. Solo el 26 de octubre tiene 1\,359\,824 escuchas, el 28~\% del archivo, y el 27 de octubre tiene apenas 4 porque el volcado se cortó a las 00:01. Este detalle reaparece al interpretar los KPIs.

\subsection{Validación}

El script \code{scripts/validate_dataset.py} recorre el archivo en streaming, sin cargarlo en memoria, y comprueba cada línea. Un registro es válido si es un objeto JSON con \code{user_name} no vacío, \code{timestamp} entero, \code{track_metadata} como objeto y nombre de artista y de canción no vacíos. Además cuenta por separado los registros sin \code{user_id} entero y sin \code{recording_msid}, registra las fechas mínima y máxima y guarda los primeros errores para mostrarlos.

El script termina con error si hay algún registro inválido, si el archivo está vacío o si la cantidad de registros no es exactamente la esperada para este volcado (4\,855\,809). La tabla~\ref{tab:validacion} resume sus opciones.

\begin{table}[H]
\centering
\caption{Argumentos de \code{validate_dataset.py}}
\label{tab:validacion}
\begin{tabularx}{\textwidth}{l l Y}
\toprule
Argumento & Por defecto & Uso \\
\midrule
\code{--file} & \code{data/processed/listens_full.jsonl} & Archivo a validar \\
\code{--expected-count} & 4855809 & Cantidad exigida; 0 desactiva el control \\
\code{--limit} & 0 & Valida solo las primeras N líneas, sin exigir la cantidad total \\
\bottomrule
\end{tabularx}
\end{table}

Con el archivo completo la validación no encontró registros inválidos, y la preparación de Neo4j comprobó además que las 4\,855\,809 filas tienen \code{user_id} entero y \code{recording_msid} no vacío. La detección de duplicados la hace cada loader al cargar, como se explica en las secciones de cada base.
